% Mira 0.11 contest-final XeLaTeX template.
% Canonical project location: paper/main.tex
\documentclass[UTF8,zihao=-4,a4paper]{ctexart}

\usepackage[left=2.5cm,right=2.5cm,top=2.6cm,bottom=2.4cm]{geometry}
\usepackage{amsmath,amssymb,mathtools,bm}
\usepackage{booktabs,array,tabularx,threeparttable}
\usepackage{siunitx}
\usepackage{graphicx,subcaption}
\usepackage{caption}
\usepackage{float}
\usepackage{fancyhdr}
\usepackage{titlesec}
\usepackage[hidelinks]{hyperref}
\usepackage{cleveref}
\usepackage{xcolor}

\definecolor{MiraInk}{RGB}{0,0,0}
\definecolor{MiraAccent}{RGB}{43,92,116}
\sisetup{detect-all,table-number-alignment=center,round-mode=places,round-precision=3}
\captionsetup{font=small,labelsep=quad,justification=centering}
\titleformat{\section}{\color{black}\heiti\zihao{3}}{\thesection}{0.8em}{}
\titleformat{\subsection}{\color{black}\heiti\zihao{4}}{\thesubsection}{0.7em}{}
\titleformat{\subsubsection}{\color{black}\heiti\zihao{-4}}{\thesubsubsection}{0.6em}{}
\pagestyle{fancy}
\fancyhf{}
\fancyhead[C]{\small 全国大学生数学建模竞赛论文}
\fancyfoot[C]{\thepage}
\renewcommand{\headrulewidth}{0.4pt}

\newcommand{\MiraTitle}[1]{%
  \begin{center}
    {\color{black}\heiti\zihao{2}\bfseries #1\par}
  \end{center}
}
\newcommand{\keywords}[1]{\par\noindent\textbf{关键词：}#1}

\begin{document}
\pagenumbering{Roman}
\MiraTitle{论文题目}

\begin{abstract}
\textbf{问题一：}解释任务的数学实质，说明方法、关键结果及直接验证证据。

\textbf{问题二：}说明共享数学核心与本问新增约束，报告关键数值、误差和稳健性证据。

\textbf{问题三：}概括算法链和主要结论，并说明适用边界、局限或可迁移条件。

\textbf{问题四：}给出范围认证、最优性证据与最终结论，避免仅凭粗网格宣称全局最优。

\keywords{数学建模；关键词二；关键词三；稳健性分析}
\end{abstract}
\clearpage

\tableofcontents
\clearpage
\pagenumbering{arabic}
\label{mira-body-start}

\section{问题重述与分析}
在此明确输入、输出、约束、评价指标与各问题之间的依赖关系。正文中的变量均写入数学环境，例如 $x_i$、$v(t)$ 与 $\rho_{\max}$。

\section{模型假设与符号说明}
假设应说明必要性、可能偏差和适用范围。正式表格统一使用三线表，如\cref{tab:symbols} 所示。

\begin{table}[H]
  \centering
  \caption{主要符号与单位}
  \label{tab:symbols}
  \begin{tabular}{cll}
    \toprule
    符号 & 含义 & 单位 \\
    \midrule
    $x_i$ & 第 $i$ 个状态变量 & \si{m} \\
    $v(t)$ & 时刻 $t$ 的速度 & \si{m.s^{-1}} \\
    \bottomrule
  \end{tabular}
\end{table}

\section{模型建立与求解}
给出变量定义、目标函数、约束条件、求解域来源和算法流程。关键公式应编号并解释，例如
\begin{equation}
  \min_{x\in\Omega} f(x) \quad \text{s.t.}\quad g_j(x)\leq 0,\;j=1,\ldots,m.
  \label{eq:optimization}
\end{equation}
式\eqref{eq:optimization} 中，$\Omega$ 的边界必须由题意或可验证的物理条件给出。

\section{结果、验证与图文论证}
每张主图之前先说明读者需要观察的问题，图后用独立段落给出可由图直接支持的结论。统计、优化或几何方法的来源应在正文真实引用\cite{sample2024}。

\section{误差分析、模型评价与推广}
误差分析至少覆盖误差来源、残差或不确定性量化、敏感性或稳健性实验、图表证据以及对结论的影响。推广部分说明可迁移数学核心、需要修改的假设、适用边界和新场景验证方式。

\label{mira-body-end}
\clearpage
\bibliographystyle{gbt7714-numerical}
\bibliography{references}

\end{document}
